% ═══════════════════════════════════════════════════════════════════════════════
%  GeoSpatial Site Readiness Analyzer — Project Documentation (LaTeX)
%  AI-Powered Location Intelligence for Commercial Real Estate & Infrastructure
% ═══════════════════════════════════════════════════════════════════════════════

\documentclass[12pt, a4paper]{article}

% ─── Packages ────────────────────────────────────────────────────────────────
\usepackage[utf8]{inputenc}
\usepackage[T1]{fontenc}
\usepackage{lmodern}
\usepackage[margin=2cm]{geometry}
\usepackage{graphicx}
\usepackage{xcolor}
\usepackage{hyperref}
\usepackage{enumitem}
\usepackage{booktabs}
\usepackage{tabularx}
\usepackage{longtable}
\usepackage{fancyhdr}
\usepackage{titlesec}
\usepackage{tcolorbox}
\usepackage{amsmath}
\usepackage{amssymb}
\usepackage{listings}
\usepackage{tikz}
\usepackage{float}
\usepackage{caption}
\usepackage{setspace}
\usepackage{parskip}

\tcbuselibrary{skins, breakable}

% ─── Color Definitions ──────────────────────────────────────────────────────
\definecolor{navy}{HTML}{0f172a}
\definecolor{darkblue}{HTML}{1e3a5f}
\definecolor{blue}{HTML}{3b82f6}
\definecolor{lightblue}{HTML}{dbeafe}
\definecolor{cyan}{HTML}{06b6d4}
\definecolor{emerald}{HTML}{10b981}
\definecolor{amber}{HTML}{f59e0b}
\definecolor{rose}{HTML}{ef4444}
\definecolor{purple}{HTML}{8b5cf6}
\definecolor{gray}{HTML}{64748b}
\definecolor{lightgray}{HTML}{f1f5f9}
\definecolor{codebg}{HTML}{f8fafc}
\definecolor{codeborder}{HTML}{e2e8f0}

% ─── Hyperref Setup ─────────────────────────────────────────────────────────
\hypersetup{
    colorlinks=true,
    linkcolor=darkblue,
    urlcolor=blue,
    citecolor=purple,
    pdftitle={GeoSpatial Site Readiness Analyzer — Documentation},
    pdfauthor={Team},
    pdfsubject={AI-Powered Location Intelligence},
}

% ─── Section Formatting ─────────────────────────────────────────────────────
\titleformat{\section}
  {\LARGE\bfseries\color{navy}}
  {\thesection.}{0.5em}{}
  [\vspace{-0.5em}\textcolor{codeborder}{\rule{\textwidth}{0.5pt}}]

\titleformat{\subsection}
  {\large\bfseries\color{darkblue}}
  {\thesubsection}{0.5em}{}

\titleformat{\subsubsection}
  {\normalsize\bfseries\color{blue}}
  {\thesubsubsection}{0.5em}{}

% ─── Header / Footer ────────────────────────────────────────────────────────
\pagestyle{fancy}
\fancyhf{}
\fancyhead[L]{\small\textcolor{gray}{GeoSpatial Site Readiness Analyzer}}
\fancyhead[R]{\small\textcolor{gray}{Project Documentation}}
\fancyfoot[C]{\small\textcolor{gray}{\thepage}}
\renewcommand{\headrulewidth}{0.4pt}
\renewcommand{\footrulewidth}{0pt}

% ─── Code Listing Style ─────────────────────────────────────────────────────
\lstdefinestyle{codestyle}{
    backgroundcolor=\color{codebg},
    basicstyle=\ttfamily\small,
    breaklines=true,
    frame=single,
    rulecolor=\color{codeborder},
    framesep=6pt,
    xleftmargin=8pt,
    xrightmargin=8pt,
    tabsize=2,
    showstringspaces=false,
    keywordstyle=\color{purple}\bfseries,
    commentstyle=\color{gray}\itshape,
    stringstyle=\color{emerald},
}
\lstset{style=codestyle}

% ─── Custom Boxes ────────────────────────────────────────────────────────────
\newtcolorbox{conceptbox}[1][]{
    colback=lightblue,
    colframe=blue,
    fonttitle=\bfseries\color{darkblue},
    title={#1},
    boxrule=0.8pt,
    arc=4pt,
    left=8pt, right=8pt, top=6pt, bottom=6pt,
    breakable,
}

\newtcolorbox{notebox}{
    colback=lightblue,
    colframe=blue,
    boxrule=0.8pt,
    arc=4pt,
    left=8pt, right=8pt, top=6pt, bottom=6pt,
    breakable,
}

\newtcolorbox{warnbox}{
    colback=rose!10,
    colframe=rose,
    boxrule=0.8pt,
    arc=4pt,
    left=8pt, right=8pt, top=6pt, bottom=6pt,
}

% ─── Custom Commands ────────────────────────────────────────────────────────
\newcommand{\tech}[1]{\texttt{\textcolor{purple}{#1}}}
\newcommand{\concept}[1]{\textbf{\textcolor{darkblue}{#1}}}
\newcommand{\score}[1]{\texttt{\textcolor{emerald}{#1}}}

% ═══════════════════════════════════════════════════════════════════════════════
%  DOCUMENT START
% ═══════════════════════════════════════════════════════════════════════════════
\begin{document}

% ─── Title Page ──────────────────────────────────────────────────────────────
\begin{titlepage}
    \centering
    \vspace*{3cm}

    {\Huge\bfseries\textcolor{navy}{GeoSpatial Site Readiness Analyzer}}

    \vspace{0.8cm}
    {\Large\textcolor{gray}{AI-Powered Location Intelligence for\\Commercial Real Estate \& Infrastructure}}

    \vspace{1.5cm}
    {\large\textcolor{gray}{Gujarat, India — Project Documentation}}

    \vspace{2cm}

    \begin{notebox}
        \textbf{What is this project?}\\[4pt]
        A full-stack web application that evaluates any location in Gujarat for commercial
        viability by analyzing 5 geospatial data layers — demographics, transportation,
        points of interest, land use, and environmental risk — and computing a composite
        site readiness score (0--100) with interactive map visualizations.
    \end{notebox}

    \vspace{1.5cm}

    \begin{tabular}{@{} l l l @{}}
        \toprule
        \textbf{Component} & \textbf{Technology} & \textbf{Purpose} \\
        \midrule
        Backend        & Python, FastAPI, Uvicorn      & REST API, scoring engine \\
        Frontend       & React 19, Vite 6, Leaflet.js  & Interactive map UI \\
        Spatial        & H3 (Uber), DBSCAN, Haversine  & Hexagonal binning, clustering \\
        Data Format    & GeoJSON                       & Geospatial feature encoding \\
        PDF Export     & ReportLab                     & Professional reports \\
        Routing        & OSRM (OpenStreetMap)          & Drive-time isochrones \\
        \bottomrule
    \end{tabular}

    \vfill
    {\small\textcolor{gray}{\today}}
\end{titlepage}

% ─── Table of Contents ──────────────────────────────────────────────────────
\tableofcontents
\newpage

% ═══════════════════════════════════════════════════════════════════════════════
%  SECTION 1: PROBLEM STATEMENT
% ═══════════════════════════════════════════════════════════════════════════════
\section{Problem Statement \& Motivation}

Selecting optimal locations for new retail stores, warehouses, EV charging stations,
telecom towers, or renewable energy installations requires analyzing dozens of
geospatial factors: population density, road network accessibility, competitor proximity,
land use regulations, environmental risks, and utility infrastructure availability.

Currently, this analysis is performed \textbf{manually} by teams of analysts using
disconnected GIS tools, spreadsheets, and gut instinct, leading to \textbf{months-long
evaluation cycles} and suboptimal site selections.

Our solution automates this entire process into an interactive, AI-powered web
application that can score any location in seconds, visualize results on an
interactive map, and generate professional PDF reports for stakeholders.

\subsection{Key Capabilities}

\begin{itemize}[leftmargin=2em]
    \item \textbf{Multi-Layer Analysis}: Ingest 5 distinct geospatial data layers simultaneously
    \item \textbf{Composite Scoring}: Weighted, explainable 0--100 score with grade (A+ to F)
    \item \textbf{Spatial Clustering}: Automatic hot-spot and cold-spot detection using DBSCAN
    \item \textbf{H3 Heatmap}: Uber's hexagonal spatial indexing for consistent area comparison
    \item \textbf{Isochrone Analysis}: Drive-time catchment areas with population estimation
    \item \textbf{Site Comparison}: Side-by-side ranking of multiple candidate sites
    \item \textbf{PDF Export}: Professional report generation with ReportLab
\end{itemize}

\newpage

% ═══════════════════════════════════════════════════════════════════════════════
%  SECTION 2: SYSTEM ARCHITECTURE
% ═══════════════════════════════════════════════════════════════════════════════
\section{System Architecture}

The application follows a \concept{client-server architecture} with a clear separation
between the React frontend and the FastAPI backend, communicating via REST API calls
proxied through the Vite development server.

\subsection{Architecture Components}

\begin{table}[H]
\centering
\begin{tabularx}{\textwidth}{l l X}
    \toprule
    \textbf{Layer} & \textbf{Technology} & \textbf{Responsibility} \\
    \midrule
    Presentation    & React 19 + Leaflet.js   & Interactive map, sidebar controls, score panels \\
    API Gateway     & Vite Dev Proxy          & Routes \texttt{/api/*} requests to backend (port 8000) \\
    Application     & FastAPI (Python)        & Business logic, scoring, clustering, export \\
    Data Access     & GeoPandas + In-Memory   & Spatial queries on GeoJSON datasets \\
    Spatial Index   & H3 Hexagonal Grid       & Uniform area partitioning for heatmaps \\
    External Svc    & OSRM API                & Real road-network routing for isochrones \\
    \bottomrule
\end{tabularx}
\caption{Architecture layer breakdown}
\end{table}

\subsection{Data Flow}

When a user clicks on the map to score a location, the following sequence occurs:

\begin{enumerate}[leftmargin=2em]
    \item \textbf{User clicks on map} $\rightarrow$ Leaflet captures \texttt{(lat, lng)} coordinates
    \item \textbf{Frontend sends} \texttt{POST /api/score} $\rightarrow$ \texttt{\{lat, lng, weights, preset, radius\_km\}}
    \item \textbf{Vite proxy forwards} to \texttt{http://localhost:8000/api/score}
    \item \textbf{Scoring Engine runs 5 sub-scorers} $\rightarrow$ each queries nearby GeoJSON features
    \item \textbf{Composite score computed} $\rightarrow$ weighted sum with threshold violation checks
    \item \textbf{JSON response returned} $\rightarrow$ \texttt{\{composite\_score, grade, sub\_scores[], violations[]\}}
    \item \textbf{ScorePanel renders} $\rightarrow$ SVG gauge, bar charts, layer-specific details
\end{enumerate}

\subsection{File Structure}

\begin{table}[H]
\centering
\small
\begin{tabularx}{\textwidth}{l l X}
    \toprule
    \textbf{Directory} & \textbf{Key Files} & \textbf{Purpose} \\
    \midrule
    \texttt{backend/}          & \texttt{main.py}           & FastAPI app with 12 REST endpoints \\
    \texttt{backend/}          & \texttt{config.py}         & Gujarat cities, weights, scoring thresholds \\
    \texttt{backend/data/}     & \texttt{generate\_data.py} & Synthetic GeoJSON generator for Gujarat \\
    \texttt{backend/services/} & \texttt{scoring.py}        & 5 sub-scorers + composite score engine \\
    \texttt{backend/services/} & \texttt{clustering.py}     & H3 heatmap + DBSCAN clustering \\
    \texttt{backend/services/} & \texttt{routing.py}        & OSRM isochrone generation \\
    \texttt{backend/services/} & \texttt{data\_ingestion.py} & GeoJSON loading + radius queries \\
    \texttt{backend/utils/}    & \texttt{spatial.py}        & Haversine distance, decay functions \\
    \texttt{frontend/src/}     & \texttt{App.jsx}           & Main state orchestrator \\
    \texttt{frontend/src/components/} & \texttt{MapView.jsx}  & Leaflet map with 5 layer renderers \\
    \texttt{frontend/src/components/} & \texttt{ScorePanel.jsx} & SVG gauge + sub-score breakdown \\
    \texttt{frontend/src/components/} & \texttt{Sidebar.jsx}    & Layer toggles, weight sliders \\
    \texttt{frontend/src/components/} & \texttt{ComparePanel.jsx} & Side-by-side site comparison \\
    \texttt{frontend/src/services/}   & \texttt{api.js}          & Axios HTTP client for all endpoints \\
    \bottomrule
\end{tabularx}
\caption{Project file structure}
\end{table}

\newpage

% ═══════════════════════════════════════════════════════════════════════════════
%  SECTION 3: DATA LAYER
% ═══════════════════════════════════════════════════════════════════════════════
\section{Data Layer — Synthetic GeoJSON Generation}

The system uses \textbf{synthetic but realistic} GeoJSON data generated for Gujarat, India.
The data generator creates features around 12 major Gujarat cities (Ahmedabad, Surat,
Vadodara, Rajkot, Gandhinagar, Bhavnagar, Jamnagar, Junagadh, Anand, Nadiad, Mehsana,
and Bharuch) with properties that mirror real-world distributions.

\subsection{Demographics Layer}

Contains census-tract-like polygons with population data. Each feature is a circular
polygon (approximated) representing a neighborhood with properties including:
\textbf{population} (5,000--500,000), \textbf{population\_density} (per km\textsuperscript{2}),
\textbf{median\_income} (\rupee 8,000--\rupee 80,000/month), \textbf{median\_age}, and
\textbf{literacy\_rate}.

\begin{conceptbox}[Concept: GeoJSON Polygon]
    A closed shape defined by an array of \texttt{[longitude, latitude]} coordinate pairs.
    The first and last coordinate must be identical to close the ring. Polygons can have
    holes (inner rings) and are used to represent areas like census tracts, zones, and
    flood regions.
\end{conceptbox}

\subsection{Transportation Layer}

Includes two geometry types: \textbf{LineString} features for roads (national highways,
state highways, arterials) with properties like \textit{lanes} and \textit{surface\_quality},
and \textbf{Point} features for transit stops (bus stations, railway stations, metro)
with \textit{daily\_ridership} data.

\subsection{Points of Interest (POI) Layer}

Point features representing businesses categorized as:
\begin{itemize}[leftmargin=2em]
    \item \textbf{Competitor}: Rival stores that compete for the same customers
    \item \textbf{Anchor}: Malls and large stores that drive foot traffic to the area
    \item \textbf{Complementary}: Businesses that benefit each other (restaurants near retail)
\end{itemize}
Each has a \textit{rating} (1--5) and \textit{city} attribute.

\subsection{Land Use \& Zoning Layer}

Polygon features representing zoned areas: \textbf{commercial}, \textbf{residential},
\textbf{industrial}, \textbf{mixed\_use}, \textbf{agricultural}, \textbf{institutional}.
Properties include \textit{development\_status} (developed/developing/undeveloped),
\textit{building\_coverage\_ratio}, and Gujarat-specific \textit{is\_gidc} flag for
Gujarat Industrial Development Corporation estates.

\subsection{Environmental Risk Layer}

Mixed geometry: \textbf{Polygon} features for flood zones (coastal Gujarat, river basins)
and earthquake risk zones (Kutch region), plus \textbf{Point} features for air quality
monitoring stations with \textit{aqi\_value} (Air Quality Index) and \textit{severity} levels.

\begin{conceptbox}[Concept: GeoJSON Standard (RFC 7946)]
    An open standard format for encoding geographic data structures using JSON.
    It supports geometry types: Point, LineString, Polygon, MultiPoint, MultiLineString,
    MultiPolygon, and GeometryCollection. Each Feature has a \texttt{geometry} and
    arbitrary \texttt{properties}.
\end{conceptbox}

\newpage

% ═══════════════════════════════════════════════════════════════════════════════
%  SECTION 4: SCORING ENGINE
% ═══════════════════════════════════════════════════════════════════════════════
\section{Scoring Engine — Multi-Criteria Decision Analysis}

The core of the application is the \concept{Scoring Engine} — a configurable system that
evaluates any location by computing 5 independent sub-scores and combining them into a
single composite score using a \concept{Weighted Linear Combination (WLC)} approach.

\subsection{Scoring Workflow}

\begin{enumerate}[leftmargin=2em]
    \item Receive target coordinates \texttt{(lat, lng)} and analysis radius (default 5\,km)
    \item For each of the 5 layers, run a spatial radius query to find nearby features
    \item Compute a raw sub-score (0--100) for each layer using domain-specific logic
    \item Apply distance decay: closer features have more influence than farther ones
    \item Multiply each sub-score by its configurable weight (weights sum to 1.0)
    \item Sum weighted scores to get the composite score
    \item Check threshold constraints (min population, max flood risk, etc.)
    \item Apply penalty for each threshold violation ($-5$ points each)
    \item Assign a letter grade (A+ through F) based on final score
\end{enumerate}

\subsection{Composite Score Formula}

The composite score $S$ is computed as:

\begin{equation}
    S = \max\left(0, \sum_{i=1}^{5} w_i \cdot s_i - 5 \cdot V\right)
\end{equation}

where:
\begin{itemize}[leftmargin=2em]
    \item $w_i$ = weight for layer $i$ (sums to 1.0)
    \item $s_i$ = sub-score for layer $i$ (0--100)
    \item $V$ = number of threshold violations
\end{itemize}

\subsection{Default Weight Configuration}

\begin{table}[H]
\centering
\begin{tabularx}{\textwidth}{l c X}
    \toprule
    \textbf{Layer} & \textbf{Weight} & \textbf{Rationale} \\
    \midrule
    Demographics      & 25\% & Population density and income are primary demand indicators \\
    Transportation    & 20\% & Highway/transit access affects customer reach \\
    Points of Interest & 20\% & Competitive landscape and anchor store proximity \\
    Land Use \& Zoning & 15\% & Regulatory compliance and development readiness \\
    Environmental Risk & 20\% & Flood/earthquake/air quality affect long-term viability \\
    \bottomrule
\end{tabularx}
\caption{Default scoring weights}
\end{table}

\subsection{Use-Case Presets}

Different business types prioritize different factors. The system supports
\textbf{weight presets} that automatically adjust scoring priorities:

\begin{table}[H]
\centering
\small
\begin{tabular}{l c c c c c l}
    \toprule
    \textbf{Use Case} & \textbf{Demo} & \textbf{Trans} & \textbf{POI} & \textbf{Land} & \textbf{Env} & \textbf{Key Priority} \\
    \midrule
    Retail Store  & 30\% & 20\% & 25\% & 10\% & 15\% & Population \& competitors \\
    EV Charging   & 15\% & 35\% & 15\% & 15\% & 20\% & Highway access \\
    Warehouse     & 10\% & 30\% & 10\% & 30\% & 20\% & Land \& roads \\
    Telecom Tower & 25\% & 15\% & 10\% & 20\% & 30\% & Coverage \& environment \\
    Solar Farm    & 5\%  & 15\% & 5\%  & 30\% & 45\% & Land \& low risk \\
    \bottomrule
\end{tabular}
\caption{Weight presets for different use cases}
\end{table}

\subsection{Sub-Score Computation Details}

\subsubsection{Demographics Sub-Score}

Queries census tracts within the radius and computes three components:

\begin{itemize}[leftmargin=2em]
    \item \textbf{Population Score (40\%)}: $\text{normalize}(\text{total\_population}, 0, 500{,}000)$
    \item \textbf{Density Score (30\%)}: $\text{normalize}(\text{avg\_density}, 0, 30{,}000/\text{km}^2)$
    \item \textbf{Income Score (30\%)}: $\text{normalize}(\text{avg\_income}, 8{,}000, 80{,}000)$
\end{itemize}

Each feature's contribution is weighted by an \concept{exponential distance decay} function.

\subsubsection{Transportation Sub-Score}

\begin{itemize}[leftmargin=2em]
    \item \textbf{Highway Score (35\%)}: Exponential decay from nearest national/state highway
    \item \textbf{Transit Score (30\%)}: Number of bus/rail/metro stops within radius
    \item \textbf{Road Density (20\%)}: Count of road features within radius
    \item \textbf{Ridership Score (15\%)}: Average daily ridership of nearby transit stops
\end{itemize}

\subsubsection{Points of Interest Sub-Score}

\begin{itemize}[leftmargin=2em]
    \item \textbf{Competitive Density (35\%)}: Uses a bell-curve formula — some competition is
          good (indicates demand), but too much saturates the market
    \item \textbf{Anchor Proximity (30\%)}: Exponential decay from nearest anchor (mall, major store)
    \item \textbf{Complementary Density (35\%)}: Count of complementary businesses nearby
\end{itemize}

\subsubsection{Land Use \& Zoning Sub-Score}

\begin{itemize}[leftmargin=2em]
    \item \textbf{Zone Type Score (40\%)}: commercial$=$95, mixed\_use$=$85, residential$=$55, industrial$=$30
    \item \textbf{Development Status (25\%)}: developing$=$90, developed$=$80, undeveloped$=$50
    \item \textbf{Building Coverage (20\%)}: Optimal around 0.4 (room for new construction)
    \item \textbf{Zone Diversity (15\%)}: More zone types nearby $=$ more versatile area
\end{itemize}

\subsubsection{Environmental Risk Sub-Score (Inverted)}

This score is \textbf{inverted} — lower risk means higher score:

\begin{itemize}[leftmargin=2em]
    \item \textbf{Flood Risk (35\%)}: Based on severity of nearest flood zone $\times$ distance decay
    \item \textbf{Earthquake Risk (35\%)}: Based on severity of nearest seismic zone $\times$ decay
    \item \textbf{Air Quality (30\%)}: $\text{normalize}(300 - \text{avg\_AQI}, 0, 300)$ — lower AQI is better
\end{itemize}

\begin{conceptbox}[Concept: Weighted Linear Combination (WLC)]
    A Multi-Criteria Decision Analysis (MCDA) method where the overall score is
    computed as:
    \[
        S = \sum_{i=1}^{n} w_i \cdot s_i
    \]
    where $w_i$ is the weight and $s_i$ is the sub-score for layer $i$.
    Weights must satisfy $\sum w_i = 1.0$. WLC is widely used in GIS-based
    suitability analysis because it is transparent, configurable, and easy to explain.
\end{conceptbox}

\subsection{Grading Scale}

\begin{table}[H]
\centering
\begin{tabular}{c c l}
    \toprule
    \textbf{Score Range} & \textbf{Grade} & \textbf{Interpretation} \\
    \midrule
    90--100 & A+ & Exceptional — top-tier location \\
    80--89  & A  & Excellent — strong across all factors \\
    70--79  & B+ & Very Good — minor weaknesses \\
    60--69  & B  & Good — suitable with some trade-offs \\
    50--59  & C+ & Average — noticeable gaps \\
    40--49  & C  & Below Average — significant concerns \\
    30--39  & D  & Poor — major issues in multiple areas \\
    0--29   & F  & Failing — unsuitable for the use case \\
    \bottomrule
\end{tabular}
\caption{Letter grade interpretation}
\end{table}

\newpage

% ═══════════════════════════════════════════════════════════════════════════════
%  SECTION 5: SPATIAL ALGORITHMS
% ═══════════════════════════════════════════════════════════════════════════════
\section{Spatial Algorithms — H3, DBSCAN, Haversine}

\subsection{Haversine Distance Formula}

Computes the \concept{great-circle distance} between two points on Earth's surface
given their latitude and longitude. This is more accurate than Euclidean distance
because it accounts for Earth's curvature.

\begin{equation}
    a = \sin^2\!\left(\frac{\Delta\phi}{2}\right) + \cos(\phi_1) \cdot \cos(\phi_2) \cdot \sin^2\!\left(\frac{\Delta\lambda}{2}\right)
\end{equation}

\begin{equation}
    c = 2 \cdot \text{atan2}\!\left(\sqrt{a},\; \sqrt{1-a}\right)
\end{equation}

\begin{equation}
    d = R \cdot c \qquad \text{where } R = 6371 \text{ km (Earth's mean radius)}
\end{equation}

Here $\phi$ represents latitude, $\lambda$ represents longitude, and $\Delta$ denotes
the difference between two points. The result $d$ is the shortest distance along
the Earth's surface.

\textbf{Used in:} Spatial radius queries, distance decay computation, DBSCAN neighbor search.

\subsection{Exponential Distance Decay}

Models the idea that closer features have more influence on a location's score.
The decay function is:

\begin{equation}
    \text{decay}(d) = e^{-\lambda d}
\end{equation}

where $d$ is distance in km and $\lambda$ is the decay rate (typically 0.2--0.4).

\begin{table}[H]
\centering
\begin{tabular}{c c c}
    \toprule
    \textbf{Distance (km)} & $\boldsymbol{\lambda = 0.3}$ & \textbf{Interpretation} \\
    \midrule
    0   & 1.000 & Full influence \\
    1   & 0.741 & 74\% influence \\
    2   & 0.549 & 55\% influence \\
    5   & 0.223 & 22\% influence \\
    10  & 0.050 & 5\% influence \\
    \bottomrule
\end{tabular}
\caption{Exponential decay values at $\lambda = 0.3$}
\end{table}

\subsection{H3 Hexagonal Spatial Index (Uber)}

\concept{H3} is a hierarchical geospatial indexing system developed by Uber. It
partitions the Earth's surface into hexagonal cells at multiple resolutions.

\subsubsection{Why Hexagons Over Squares?}

\begin{itemize}[leftmargin=2em]
    \item Every neighbor is equidistant (6 neighbors, all same distance from center)
    \item No orientation bias (squares have diagonal vs.\ cardinal neighbor asymmetry)
    \item Better approximation of circular search areas
    \item Hierarchical: each parent hex contains $\approx$7 child hexes
\end{itemize}

\subsubsection{Resolution Levels}

\begin{table}[H]
\centering
\begin{tabularx}{\textwidth}{c c c X}
    \toprule
    \textbf{Resolution} & \textbf{Hex Area} & \textbf{Edge Length} & \textbf{Use in This Project} \\
    \midrule
    4 & $\approx$1,770 km\textsuperscript{2} & $\approx$22.6 km & State-level overview \\
    5 & $\approx$253 km\textsuperscript{2}   & $\approx$8.5 km  & Regional analysis \\
    7 & $\approx$5.16 km\textsuperscript{2}  & $\approx$1.22 km & City-level heatmaps (default) \\
    9 & $\approx$0.105 km\textsuperscript{2} & $\approx$174 m   & Street-level analysis \\
    \bottomrule
\end{tabularx}
\caption{H3 resolution levels and their use cases}
\end{table}

\subsubsection{Key H3 Functions (v4 API)}

\begin{lstlisting}[language=Python]
import h3

# Convert lat/lng to hex cell ID
hex_id = h3.latlng_to_cell(lat, lng, resolution)

# Get center coordinates of a hex
center_lat, center_lng = h3.cell_to_latlng(hex_id)

# Get polygon boundary vertices
boundary = h3.cell_to_boundary(hex_id)
\end{lstlisting}

\subsection{DBSCAN Clustering Algorithm}

\concept{DBSCAN (Density-Based Spatial Clustering of Applications with Noise)}
groups nearby high-scoring hexagons into ``hot spots'' and low-scoring hexagons
into ``cold spots'' (underserved areas).

\subsubsection{Advantages Over K-Means}

\begin{itemize}[leftmargin=2em]
    \item Does \textbf{not} require specifying the number of clusters in advance
    \item Can find clusters of \textbf{arbitrary shape} (not just spherical)
    \item Identifies \textbf{noise points} (outliers that don't belong to any cluster)
    \item Uses two parameters: \textbf{eps} (max neighbor distance) and \textbf{min\_samples} (min cluster size)
\end{itemize}

\subsubsection{Algorithm Steps}

\begin{enumerate}[leftmargin=2em]
    \item For each unvisited hexagon, find all hexagons within $\epsilon$ km distance (using Haversine)
    \item If $\geq$ \texttt{min\_samples} neighbors found, start a new cluster
    \item \textbf{Expand}: add each neighbor's neighbors if they also meet the \texttt{min\_samples} threshold
    \item Continue until no more points can be added to the cluster
    \item Unvisited points remaining are classified as \textbf{noise}
\end{enumerate}

\subsubsection{Hot Spot vs.\ Cold Spot Detection}

The system runs DBSCAN twice:
\begin{itemize}[leftmargin=2em]
    \item \textbf{Hot spots}: Cluster hexagons with scores $> 60$ (high potential areas)
    \item \textbf{Cold spots}: Cluster hexagons with scores $< 35$ (underserved/risky areas)
\end{itemize}

\subsection{Competitive Density Score (Bell Curve)}

For retail site selection, some competition is desirable (validates demand),
but too many competitors saturate the market:

\begin{equation}
    \text{score}(n) =
    \begin{cases}
        n / n_{\min} & \text{if } n < n_{\min} \\
        1.0 & \text{if } n_{\min} \leq n \leq n_{\max} \\
        \max\!\left(0,\; 1 - \dfrac{n - n_{\max}}{2 \cdot n_{\max}}\right) & \text{if } n > n_{\max}
    \end{cases}
\end{equation}

where $n$ is the number of competitors, $n_{\min} = 2$, $n_{\max} = 6$ (default).

\newpage

% ═══════════════════════════════════════════════════════════════════════════════
%  SECTION 6: ROUTING & ISOCHRONE
% ═══════════════════════════════════════════════════════════════════════════════
\section{Routing \& Isochrone Analysis}

An \concept{isochrone} is a polygon on a map that shows all areas reachable within a
given travel time from a point. Unlike a simple radius circle, isochrones follow the
actual road network — areas with highways are reachable faster than areas with only
narrow roads.

\subsection{How Isochrones Are Computed}

\begin{enumerate}[leftmargin=2em]
    \item User clicks ``Analyze'' in the drive-time catchment section
    \item Frontend sends \texttt{POST /api/isochrone} with \texttt{lat, lng, mode, intervals}
    \item Backend sends requests to the OSRM (Open Source Routing Machine) public API
    \item OSRM returns route geometries for sample directions from the origin
    \item Backend approximates the isochrone polygon from route endpoints
    \item For each ring, population within the catchment area is estimated from demographics data
    \item Result: nested polygons (10-min, 20-min, 30-min rings) with population counts
\end{enumerate}

\begin{conceptbox}[Concept: OSRM (Open Source Routing Machine)]
    A high-performance routing engine that uses OpenStreetMap data to compute
    shortest/fastest paths on real road networks. It supports car, bicycle, and
    foot routing profiles. The public demo server (\texttt{router.project-osrm.org})
    is used for development; production would require a self-hosted instance.
\end{conceptbox}

\subsection{Catchment Population Estimation}

For each isochrone ring, the system queries the demographics layer for census tracts
whose centroids fall within the polygon, then sums their population. This tells
stakeholders:

\begin{quote}
    \textit{``Within a 10-minute drive, 378,000 potential customers reside.''}
\end{quote}

\newpage

% ═══════════════════════════════════════════════════════════════════════════════
%  SECTION 7: FRONTEND
% ═══════════════════════════════════════════════════════════════════════════════
\section{Frontend — React + Leaflet Interactive Map}

\subsection{Design System}

The UI uses a \textbf{premium dark theme} with glassmorphism effects (semi-transparent
backgrounds with backdrop blur). Key design tokens:

\begin{table}[H]
\centering
\begin{tabularx}{\textwidth}{l l X}
    \toprule
    \textbf{Token} & \textbf{Value} & \textbf{Usage} \\
    \midrule
    \texttt{--bg-primary}    & \texttt{\#0a0e1a}            & Main background \\
    \texttt{--bg-glass}      & \texttt{rgba(17,24,39,0.65)} & Glassmorphism panels \\
    \texttt{--accent-blue}   & \texttt{\#3b82f6}            & Primary interactive elements \\
    \texttt{--accent-emerald} & \texttt{\#34d399}           & Success / good scores \\
    \texttt{--accent-rose}   & \texttt{\#fb7185}            & Danger / poor scores \\
    \texttt{--font-main}     & Inter                        & UI text (Google Fonts) \\
    \texttt{--font-mono}     & JetBrains Mono               & Coordinates, scores \\
    \bottomrule
\end{tabularx}
\caption{CSS design tokens}
\end{table}

\subsection{Component Architecture}

\begin{table}[H]
\centering
\begin{tabularx}{\textwidth}{l l X}
    \toprule
    \textbf{Component} & \textbf{Responsibility} & \textbf{Key Features} \\
    \midrule
    \texttt{App.jsx}        & State orchestrator & Manages all state: layers, scores, compare, heatmap \\
    \texttt{Sidebar}        & Controls panel     & Layer toggles, weight sliders, radius, action buttons \\
    \texttt{MapView}        & Interactive map    & Leaflet with 5 layers, heatmap hexagons, clusters \\
    \texttt{ScorePanel}     & Score display      & Animated SVG gauge, 5 sub-score bars, violations \\
    \texttt{ComparePanel}   & Site comparison    & Ranked cards with mini bar charts \\
    \bottomrule
\end{tabularx}
\caption{React component architecture}
\end{table}

\subsection{Leaflet Map Layers}

The Leaflet map renders multiple overlapping layers:

\begin{itemize}[leftmargin=2em]
    \item \textbf{Base tiles}: CartoDB Dark Matter (dark theme compatible)
    \item \textbf{Demographics}: Blue polygons with opacity based on population density
    \item \textbf{Transportation}: Orange lines for roads, orange points for transit stops
    \item \textbf{POI}: Color-coded circles (red $=$ competitor, yellow $=$ anchor, green $=$ complementary)
    \item \textbf{Land Use}: Zone polygons colored by type (yellow $=$ commercial, blue $=$ residential)
    \item \textbf{Environmental}: Risk zone polygons with dashed borders for flood areas
    \item \textbf{Heatmap}: H3 hexagon polygons colored by score (green $\to$ yellow $\to$ red)
    \item \textbf{Clusters}: Large circles for hot spots (green) and cold spots (blue)
\end{itemize}

\subsection{SVG Score Gauge}

The score panel features an animated circular gauge built with SVG:

\begin{itemize}[leftmargin=2em]
    \item Two concentric circles: background track and progress arc
    \item Progress is controlled via \texttt{stroke-dashoffset} CSS property
    \item Color transitions: \textcolor{rose}{Red (0--24)} $\to$ \textcolor{amber}{Orange (25--39)}
          $\to$ \textcolor{amber}{Yellow (40--59)} $\to$ \textcolor{emerald}{Green (60--79)}
          $\to$ \textcolor{emerald}{Bright Green (80--100)}
    \item Animated with CSS transition for smooth arc drawing on load
\end{itemize}

\newpage

% ═══════════════════════════════════════════════════════════════════════════════
%  SECTION 8: API DESIGN
% ═══════════════════════════════════════════════════════════════════════════════
\section{API Design — FastAPI REST Endpoints}

\begin{table}[H]
\centering
\begin{tabularx}{\textwidth}{l l X}
    \toprule
    \textbf{Method} & \textbf{Endpoint} & \textbf{Description} \\
    \midrule
    GET  & \texttt{/api/health}         & Server health check with layer counts \\
    GET  & \texttt{/api/layers}         & List all available layers with metadata \\
    GET  & \texttt{/api/layers/\{name\}} & Fetch GeoJSON data for a specific layer \\
    POST & \texttt{/api/score}          & Compute site readiness score for a point \\
    POST & \texttt{/api/score/polygon}  & Score all hexagons within a drawn polygon \\
    GET  & \texttt{/api/h3/heatmap}     & Generate H3 hexagonal heatmap of scores \\
    POST & \texttt{/api/cluster}        & Run DBSCAN clustering for hot/cold spots \\
    POST & \texttt{/api/isochrone}      & Compute drive-time isochrone rings \\
    POST & \texttt{/api/compare}        & Compare multiple candidate sites \\
    POST & \texttt{/api/export}         & Generate PDF or JSON report \\
    GET  & \texttt{/api/presets}        & Get available use-case weight presets \\
    GET  & \texttt{/api/config}         & Get application configuration \\
    \bottomrule
\end{tabularx}
\caption{Complete API endpoint reference}
\end{table}

\subsection{Score Request / Response Example}

\begin{lstlisting}[language=Python]
# Request: POST /api/score
{
    "lat": 23.0225, "lng": 72.5714,
    "weights": {"demographics": 0.25, "transportation": 0.20,
                "poi": 0.20, "landuse": 0.15, "environmental": 0.20},
    "preset": "retail_store",
    "radius_km": 5
}

# Response:
{
    "composite_score": 51.7,
    "grade": "C+",
    "sub_scores": [
        {"layer": "demographics", "score": 56.6, "weight": 0.25,
         "weighted_score": 14.1,
         "details": {"population_nearby": 290407, "avg_income": 75926}},
        {"layer": "transportation", "score": 14.4, "weight": 0.20,
         "weighted_score": 2.9, "details": {"roads": 0, "transit_stops": 1}},
        ...
    ],
    "threshold_violations": []
}
\end{lstlisting}

\newpage

% ═══════════════════════════════════════════════════════════════════════════════
%  SECTION 9: PDF REPORT
% ═══════════════════════════════════════════════════════════════════════════════
\section{PDF Report Generation}

The system generates professional PDF reports using Python's \tech{ReportLab} library.
When a user clicks ``Export Report (PDF),'' the backend scores all pinned sites, ranks
them, and builds a multi-page PDF with:

\begin{itemize}[leftmargin=2em]
    \item \textbf{Title page} with generation timestamp
    \item \textbf{Weights table} showing the scoring configuration used
    \item \textbf{Rankings table} with all sites sorted by composite score
    \item \textbf{Detailed breakdown} for each site: layer scores, weights, weighted contributions
    \item \textbf{Threshold violations} highlighted with warning symbols
\end{itemize}

The PDF is streamed directly to the browser as a download via FastAPI's
\texttt{StreamingResponse} with \texttt{media\_type='application/pdf'}.

\newpage

% ═══════════════════════════════════════════════════════════════════════════════
%  SECTION 10: END-TO-END WORKFLOW
% ═══════════════════════════════════════════════════════════════════════════════
\section{End-to-End Workflow}

\subsection{Scenario: Evaluating a Retail Store Location Near Ahmedabad}

\begin{enumerate}[leftmargin=2em, itemsep=8pt]
    \item \textbf{Launch}: Start backend (\texttt{uvicorn}) and frontend (\texttt{npm run dev}).
          Backend loads 1,253 synthetic GeoJSON features into memory.

    \item \textbf{Explore Data}: Toggle layer toggles to visualize demographics (blue zones),
          transportation (orange roads), POIs (colored dots), land use, and environmental
          risks across Gujarat.

    \item \textbf{Score a Site}: Click near Ahmedabad on the map. The system queries all
          features within 5\,km, applies distance decay, computes 5 sub-scores, and returns
          a composite score (e.g., 51.7 / 100, Grade C+).

    \item \textbf{Analyze Breakdown}: The Score Panel shows demographics$=$56.6 (good
          population density), transportation$=$14.4 (far from highway), POI$=$41.3
          (moderate competition), land\_use$=$58.9 (mixed-use zone), environmental$=$87.8
          (low risk).

    \item \textbf{Catchment Analysis}: Click ``Analyze'' in Drive-Time Catchment. The system
          computes 10/20/30-minute drive-time isochrones via OSRM, showing 378K people within 10 min.

    \item \textbf{Generate Heatmap}: Click ``Generate H3 Heatmap'' to see color-coded hexagons
          across Gujarat. Green hexagons indicate high-potential areas; red indicates poor candidates.

    \item \textbf{Find Hot Spots}: Click ``Detect Hot Spots'' to run DBSCAN clustering. The
          system identifies clusters of high-scoring hexagons (investment opportunities) and
          low-scoring clusters (underserved areas).

    \item \textbf{Compare Sites}: Enable compare mode, click 3--4 candidate locations. The
          Compare Panel shows ranked cards with side-by-side bar charts.

    \item \textbf{Export Report}: Click ``Export Report (PDF)'' to download a professional
          report with all site data, rankings, and score breakdowns.
\end{enumerate}

\newpage

% ═══════════════════════════════════════════════════════════════════════════════
%  SECTION 11: GLOSSARY
% ═══════════════════════════════════════════════════════════════════════════════
\section{Key Concepts Glossary}

\begin{longtable}{@{} p{4cm} p{12cm} @{}}
    \toprule
    \textbf{Term} & \textbf{Definition} \\
    \midrule
    \endhead
    GeoJSON           & Open standard format (RFC 7946) for encoding geographic features using JSON \\[4pt]
    Haversine Formula & Calculates great-circle distance between two lat/lng points on a sphere \\[4pt]
    Distance Decay    & Mathematical function where influence decreases exponentially with distance \\[4pt]
    H3                & Uber's hierarchical hexagonal spatial indexing system for the globe \\[4pt]
    DBSCAN            & Density-based clustering algorithm that finds arbitrarily-shaped clusters \\[4pt]
    Isochrone         & A polygon showing all points reachable within a given travel time \\[4pt]
    OSRM              & Open Source Routing Machine — routing engine using OpenStreetMap road data \\[4pt]
    WLC               & Weighted Linear Combination — MCDA method: $S = \sum w_i \cdot s_i$ \\[4pt]
    Leaflet           & Open-source JavaScript library for interactive maps \\[4pt]
    FastAPI           & Modern high-performance Python web framework for building APIs \\[4pt]
    Glassmorphism     & UI design trend using frosted-glass effect (blur + transparency) \\[4pt]
    GeoJSON Feature   & A geographic object with geometry (shape) and properties (attributes) \\[4pt]
    Spatial Query     & Database query filtering features based on geographic criteria (e.g., within radius) \\[4pt]
    Composite Score   & A single number (0--100) combining multiple weighted sub-scores \\[4pt]
    Threshold Constraint & Minimum/maximum values that trigger penalties if violated \\[4pt]
    Heat Map          & Visualization where areas are colored by intensity of a metric \\[4pt]
    Hot Spot          & A cluster of adjacent areas with consistently high scores \\[4pt]
    Cold Spot         & A cluster of adjacent areas with consistently low scores (underserved) \\[4pt]
    Catchment Area    & The geographic region from which a business draws its customers \\[4pt]
    CORS              & Cross-Origin Resource Sharing — browser security policy for cross-domain requests \\[4pt]
    Vite              & Next-generation frontend build tool with instant hot module replacement \\[4pt]
    ReportLab         & Python library for programmatic PDF document generation \\[4pt]
    \bottomrule
\end{longtable}

\vfill
\begin{center}
    \textcolor{gray}{\rule{5cm}{0.4pt}}\\[6pt]
    \textcolor{gray}{\textit{— End of Documentation —}}
\end{center}

\end{document}
